\documentclass[a4paper,12pt]{article}
\usepackage[T2A]{fontenc}
\usepackage[utf8]{inputenc}
\usepackage[russian]{babel}
\usepackage{amsmath,amssymb}
\usepackage{hyperref}
\usepackage{geometry}
\geometry{top=2cm, bottom=2cm, left=2.5cm, right=2.5cm}

\title{Конспект лекции 3: \\ Основные задачи Computer Vision}
\author{Липкин С.М.}
\date{}

\begin{document}
\maketitle

\section{Введение в Computer Vision}

Компьютерное зрение (CV) решает задачу преобразования визуальных данных в высокоуровневое понимание сцены. Нейронные сети стали основным инструментом в CV благодаря автоматическому извлечению признаков, иерархическому представлению (края $\to$ текстуры $\to$ объекты) и обучению по размеченным данным.

\subsection{Основные типы задач}

\begin{enumerate}
    \item \textbf{Классификация} — сеть выдаёт вероятности классов через softmax.
    \item \textbf{Детекция объектов} — локализация (bounding box) + классификация.
    \item \textbf{Сегментация}:
    \begin{itemize}
        \item Семантическая — все объекты одного класса одним цветом.
        \item Instance segmentation — каждый объект уникальным цветом.
    \end{itemize}
    \item \textbf{Распознавание лиц} — embeddings с последующим сравнением расстояний.
    \item \textbf{Детекция краев} — градиентные фильтры или CNN.
    \item \textbf{Восстановление изображений} — encoder-decoder архитектуры.
    \item \textbf{Сопоставление признаков} — keypoint extraction + matching.
    \item \textbf{Реконструкция сцены} — Structure from Motion, depth estimation.
    \item \textbf{Анализ движения} — object tracking, optical flow.
\end{enumerate}

\section{Свёрточные нейронные сети (CNN)}

\subsection{Иерархия признаков}
\begin{itemize}
    \item Layer 1: края, углы, цвета.
    \item Layer 2--3: текстуры, формы.
    \item Layer 4--5: части объектов.
    \item FC: целые объекты.
\end{itemize}

\subsection{Компоненты CNN}
\begin{enumerate}
    \item Свёрточный слой — фильтры $\to$ feature maps.
    \item ReLU: $f(x) = \max(0, x)$.
    \item Pooling — уменьшение размерности.
    \item FC слои — классификация.
\end{enumerate}

\subsection{Ключевые архитектуры}
\begin{itemize}
    \item \textbf{ResNet} — skip connections для глубоких сетей.
    \item \textbf{Inception} — параллельные свёртки разных масштабов.
    \item \textbf{EfficientNet} — композитное масштабирование (глубина, ширина, разрешение).
\end{itemize}

\section{Transfer Learning}

Предобученная сеть (backbone на ImageNet: ResNet, VGG) адаптируется под новую задачу. Стратегии обучения:
\begin{itemize}
    \item \textbf{Freeze} — фиксация ранних слоёв.
    \item \textbf{Fine-tune} — дообучение поздних слоёв.
\end{itemize}
Замена «головы» сети: классификация — новые FC слои, детекция — RPN + classifier, сегментация — decoder.

\section{Архитектуры детекции}

\subsection{Двухстадийные (Faster R-CNN)}
\begin{enumerate}
    \item RPN генерирует anchor boxes.
    \item ROI Pooling извлекает признаки.
    \item Классификатор: класс + координаты.
\end{enumerate}
Точнее, но медленнее.

\subsection{Одностадийные (YOLO, SSD)}
\begin{itemize}
    \item \textbf{YOLO}: grid-based подход, один forward pass, NMS для удаления дубликатов. Эволюция: YOLOv1 (2016) $\to$ YOLOv8 (2023), anchor-free, NAS.
    \item \textbf{SSD}: multi-scale detection, default boxes разных aspect ratios, VGG16 backbone.
\end{itemize}
Быстрее, хуже на мелких объектах.

\subsection{FPN (Feature Pyramid Network)}
Объединение высокого разрешения с высокоуровневой семантикой. Bottom-up (C2--C5) + top-down с lateral connections (P2--P5). Используется в Faster R-CNN + FPN, Mask R-CNN.

\section{Архитектуры сегментации}

\subsection{DeepLab}
Atrous (dilated) convolution увеличивает receptive field без роста параметров. ASPP (Atrous Spatial Pyramid Pooling) — параллельные свёртки с rates 1, 6, 12, 18. Encoder-decoder структура.

\subsection{Генеративные модели}
\begin{itemize}
    \item \textbf{GAN}: Generator vs Discriminator (StyleGAN, pix2pix).
    \item \textbf{Diffusion}: пошаговое удаление шума (DALL-E, Stable Diffusion).
    \item \textbf{VAE}: сжатие в latent space.
\end{itemize}

\section{Применение и тренды}

Применения: медицина, автономное вождение, промышленность, сельское хозяйство, безопасность.

Современные тренды:
\begin{itemize}
    \item Transfer Learning и предобученные модели.
    \item Vision Transformers (patches + attention).
    \item Мультимодальность (CLIP, DALL-E).
    \item Self-supervised learning.
    \item Explainable AI (Grad-CAM).
    \item Генеративные модели (GAN, Diffusion).
\end{itemize}

\end{document}
